% ICAPS 2027 submission draft v1 (2026-09-01)
% NOTE: swap to the official AAAI two-column author kit when the ICAPS-27
% kit is published; content and lengths are already calibrated for
% 8 pages + references.
\documentclass[twocolumn]{article}
\usepackage[margin=2cm]{geometry}
\usepackage{amsmath,amssymb,amsthm}
\usepackage{booktabs}
\usepackage{graphicx}
\usepackage{xcolor}
\usepackage{hyperref}
\newtheorem{definition}{Definition}
\newtheorem{proposition}{Proposition}
\newcommand{\cmax}{C_{\max}}

\title{Breakable Commitments: Unblocking Closed-Loop\\
Replanning in Coupled Job-Shop Scheduling and Multi-Agent Path Finding}
\author{Anonymous submission}
\date{}

\begin{document}
\maketitle

\begin{abstract}
Flexible job-shop scheduling (FJSP) and multi-agent path finding (MAPF)
couple naturally in smart factories: machines process operations while
AGVs move material over a shared grid. State-of-the-art stacks solve this
coupled problem \emph{once} and execute open-loop, and a growing line of
work proposes closed-loop replanning---yet we find that in a full-stack
implementation, \textbf{every} dynamic plan revision is silently
unexecutable: the request encoder cannot express successors of
in-flight commitments, so a single active transport makes a revision
impossible to activate. We make four contributions.
(1)~A formal online specification (O-IFJSP-MAPF) of the coupled problem
with \emph{commitment-consistent revision} semantics that mirrors real
execution machinery. (2)~An empirical diagnosis on a coupled benchmark:
one-shot CP-SAT plans are \emph{passively robust} to outages and duration
variance (4--14\% degradation under a 200-step machine outage), which
masks the broken revision pathway---until a job stream arrives, where
throughput collapses by two orders of magnitude.
(3)~\emph{Soft commitments}: a conservative \emph{commitment
projection} that maps each in-flight commitment $(\mu,l,r)$ to exact
resource and source $(\mu,l)$ and a bounded release $\bar r\ge r$
governed by a single interpretable penalty parameter, with a
two-level soundness guarantee---execution safety holds
\emph{unconditionally} under stack release semantics, plan-time validity
under a bounded-delay assumption. The projection lifts revision
activation from 53\% to 99\% across 2401 triggered revisions.
(4)~A 1209-episode factorial evaluation with paired significance tests
that yields a provocative conclusion: \textbf{event-triggered revision
loses to doing nothing in every disruption scenario we could
construct}---one-shot plans absorb 200-step outages within 6\% via
slack, while conservative revision pays a 25--125\% premium. Under
streaming arrivals, soft commitments deliver reliability---zero
revision failures versus $1$--$18$ per cell, guaranteed completion
versus starvation in the worst legacy cell, and $23$--$31\%$ throughput
gains on the largest shop---and a snapshot-based \emph{value-aware
trigger} converts the diagnosis into a working decision policy:
replanning only $1.3\%$ of the time, it improves on both never-replan
and always-replan in this regime ($2876$ vs.\ $3018$/$3466$ mean
makespan, paired tests). Closed-loop factory scheduling needs a
when-to-replan decision policy, not a better trigger---and a simple
one already suffices.
\end{abstract}

\section{Introduction}
Manufacturing execution increasingly couples two classical problems.
On the shop floor, \emph{flexible job-shop scheduling} assigns operations
to eligible machines and sequences them. On the logistics floor,
\emph{multi-agent path finding} routes AGVs collision-free over a grid.
The two layers interact through transport: the schedule emits a stream of
pickup-and-delivery requests whose realized travel times depend on
congestion the scheduler does not see; conversely, routing decisions
determine when material becomes available for processing.

Practical systems therefore decompose the problem hierarchically, and the
most recent integrated frameworks~\cite{li2026framework} plan once and
repair only the routing layer at execution time---a
\emph{generate-then-refine} paradigm whose authors themselves list
closing the loop as future work. Meanwhile, the replanning literature
(plan repair, execution monitoring, robust execution) offers a rich
vocabulary~\cite{fox2006plan,domshlak2013satisficing} but rarely meets
the coupled factory stack, and dynamic-scheduling studies that do model
machine breakdowns abstract transport away to a fixed time
matrix~\cite{lu2025dynamic}.

This paper asks a deliberately unglamorous question: \emph{in a realistic
coupled stack, what does it take to change the plan after execution has
started?} The answer turns out to be a mechanism-level bottleneck rather
than a policy-level one:

\begin{itemize}
\item We formalize the online coupled problem and its
  \emph{commitment-consistent revision} semantics
  (\S\ref{sec:formulation}), mirroring the ledger/activation machinery of
  a real engine (plan revisions, frozen operations, activation checks).
\item We show empirically that the standard encoding of commitment
  consistency \emph{categorically blocks} revision: any queued or
  in-transit predecessor renders the residual problem inexpressible, and
  activation is vetoed on 47\% of trigger opportunities ($n{=}2401$
  revisions)---across disruption revision, job arrivals, and periodic
  refresh (\S\ref{sec:diagnosis}).
\item We introduce \emph{soft commitments}: a bounded-release encoding
  with one interpretable penalty parameter, plus exact commitment-machine
  resolution from queue membership and transport destinations
  (\S\ref{sec:soft}). The encoding is client-side and opt-in (a
  $\sim$90-line patch in our implementation) and raises activation from
  53\% to 99\%.
\item We instantiate the $(\tau,\kappa)$ recovery-orchestrator policy
  space on top and evaluate on a 1209-episode benchmark (948-episode
  main factorial plus six sensitivity phases) spanning instance families,
  fleet sizes, disruption scenarios, and commitment modes
  (\S\ref{sec:experiments}).
\end{itemize}

Our headline findings revise two intuitions. First, \emph{one-shot plans
are passively robust}: slack plus release-by-plan semantics absorb severe
disruptions (a 200-step outage degrades makespan by only 4--14\%),
so the broken revision pathway is invisible until it is
catastrophic---on the pre-repair stack, streaming arrivals starved
entirely ($24.4\to0.24$ jobs/ksteps). Second, once revisions are
unblocked and measured at scale (1209 episodes across seven factorial
phases, paired tests at $p<10^{-4}$), event-triggered revision
\emph{loses} to doing nothing in every disruption scenario: the value
of closed-loop replanning in this stack is reliability under
arrivals---zero revision failures, guaranteed completion, and
$23$--$31\%$ throughput gains on larger shops---not disruption
response.

\section{Related Work}
\textbf{Coupled scheduling and routing.} Integrated FJSP with transport
has a long OR tradition; recent surveys catalogue metaheuristic and
exact approaches~\cite{xin2025fjsppt,dauzere2024fjsp} on classic
instance families~\cite{brandimarte1993routing}. Learning-based
dispatching transfers across
instances~\cite{zhang2020l2d,smit2025gnn}, and the closest coupled
framework co-optimizes a JSP schedule with PBS routing under a
one-shot paradigm~\cite{li2026framework}. These evaluations use single
maps and open-loop execution; revision semantics are not studied.

\textbf{MAPF and lifelong MAPF.} Optimal and bounded-suboptimal search
(CBS~\cite{sharon2015cbs}, EECBS~\cite{li2021eecbs}), scalable
suboptimal search (LaCAM~\cite{okumura2023lacam}), and learned policies
(MAPF-GPT~\cite{andreychuk2025mapfgpt}) are benchmarked on standardized
grid suites~\cite{skrynnik2024pogema}; lifelong variants handle task
streams with rolling horizons~\cite{ma2017mapd,li2021rhcr} and repair
operators~\cite{li2022lns2}. Robust MAPF treats stochastic
travel~\cite{vanderwedden2020sttcbs}. None of these model
machine eligibility, processing durations, or plan-level commitments.

\textbf{Replanning and plan repair.} Plan stability, repair, and
execution monitoring are classical concerns~\cite{fox2006plan}, with
commitment-based satisficing planning studying what to
fix~\cite{domshlak2013satisficing}. Our contribution is orthogonal and
practical: we show where a production-grade coupled stack silently
\emph{forbids} repair---the commitment encoding itself---and give the
minimal sound relaxation that unblocks it, turning the
stability--responsiveness trade-off into a single measurable penalty
parameter.

\section{Online Coupled Problem and Revision Semantics}
\label{sec:formulation}
\subsection{Static core (I-FJSP-MAPF)}
Jobs $\mathcal{J}$ carry operation chains $O_j$; each operation $o$ has
eligible machines $E(o)$ with durations $p_{o,m}$; machines occupy cells
of a grid $G$; a fleet $\mathcal{K}$ of unit-capacity AGVs moves one edge
per step. A solution assigns (i) machines and sequences, (ii) transport
tasks to AGVs, and (iii) collision-free spacetime trajectories
$\pi_k$, minimizing makespan $\cmax$. The coupling constraints are the
hand-offs: an operation may start only after its predecessor's material
\emph{actually arrives}---a time determined by the routing layer.

\subsection{Online setting (O-IFJSP-MAPF)}
An exogenous disruption process $\xi$ emits machine outages
$\texttt{m\_down}(m,d)$, AGV outages, and temporary obstacles; a job
stream releases jobs at times $a_j$. A closed-loop policy observes the
state (remaining operations, machine/AGV status, in-flight tasks, current
plan) and may, at any step, issue a \emph{revision}.

\begin{definition}[Feasible revision]
A revision $\Pi_{t+1}=\rho(\Pi_t\mid s_t)$ is feasible iff it
(i)~agrees with the executed prefix; (ii)~honors \emph{commitments}
---started operations stay on their machines, loaded AGVs deliver to
their destinations; and (iii)~is (approximately) optimal for the residual
problem.
\end{definition}

\begin{definition}[Disruption regret]
$\mathrm{RRegret}(\pi;S)=\cmax(\pi,S)/\cmax(\pi,S_0)$ for scenario $S$
against the empty scenario $S_0$.
\end{definition}

\subsection{Commitment classes}
At any instant each operation is in one of:
\textsc{completed}, \textsc{processing} (machine busy),
\textsc{machine\_queue} (queued at a machine), \textsc{transport}
(material on an AGV), or \textsc{pending} (free to schedule).
Only \textsc{pending} operations are reschedulable; the rest are frozen.
The revision machinery must therefore encode, for the first reschedulable
operation of each job, \emph{when its input material exists} (release
bound) and \emph{where} (source cell). This seemingly clerical encoding
is the paper's protagonist.

\section{Diagnosis: Commitments Block Revision}
\label{sec:diagnosis}
State-of-the-art coupled frameworks plan once and repair only the
routing layer, listing closed-loop revision as future
work~\cite{li2026framework}. We argue this is no accident but a
consequence of how such stacks naturally honor commitments. The
baseline---and, we contend, the standard---encoding of commitment
consistency in a stack that revises through a residual problem is
\emph{encode-or-veto}: express exactly what is expressible, and veto
the revision otherwise. It is a defensible design: it can never violate
a commitment. We instrument a full-stack implementation of this
baseline pattern and show what it costs. The residual-problem encoder
works as follows for a job's first reschedulable operation $o$ with
predecessor $p$:
\begin{itemize}
\item $p\in$\textsc{processing}: release $:=$ machine availability
(handled);
\item $p\in$\{\textsc{machine\_queue}, \textsc{transport}\}: the encoder
\emph{cannot express} the release or source; it appends $p$ (or $o$) to
an \texttt{unresolved} list, and the artifact is marked
\texttt{execution\_ready=False} iff that list is nonempty;
\item the solver then \emph{refuses to activate} any such revision.
\end{itemize}

Since a running factory almost always has queued or in-transit work,
\textbf{revision activation fails categorically}. In our pilots this
single pathway failure explained three seemingly unrelated phenomena:
disruption-triggered revision, job-arrival revision (starving all late
arrivals: throughput 24.4$\to$0.24 jobs/ksteps), and periodic refresh.
Simultaneously, one-shot plans looked deceptively healthy---slack absorbs
outages (4--14\% makespan growth under a 200-step outage) and duration
variance ($D\in[1.00,1.15]$)---so nothing forced anyone down the
revision path. We consider this a cautionary tale for evaluating
closed-loop machinery \emph{only} under weak disruptions.

\section{Soft Commitments as Conservative Projection}
\label{sec:soft}

\subsection{Commitment projection}
The repair is best understood not as a patch but as an abstraction.
An \emph{execution commitment} for a frozen operation $p$ is the triple
$\kappa(p)=(\mu,\,l,\,r)$: the committed resource $\mu$ (the machine
where $p$ will run, hence where its successor's input material will
land), the material source $l$ (that machine's cell), and the
earliest-safe-release $r$ (the earliest step at which the successor of
$p$ may begin). A \emph{commitment projection}
$\Pi$ maps the execution state to the residual problem instance by
turning each commitment into (i)~a release bound $\bar r\ge r$ for the
successor, (ii)~a source $\bar l$ for its transfer request, and
(iii)~machine-availability offsets.

Two projections are of interest. An \emph{exact} projection achieves
$\bar r=r$ and $\bar l=l$; a \emph{conservative} projection requires only
$\bar r\ge r$ with $\bar l=l$ exact. The legacy encoder is an
\emph{all-or-nothing} projection: any commitment whose $(\mu,l)$ it
cannot derive exactly is excluded from the encoding and vetoes
activation outright. Soft commitments complete the projection:

\paragraph{Exact $(\mu,l)$ resolution.}
Operation objects lack an assigned machine before processing starts,
but the information exists elsewhere: a \textsc{machine\_queue}
operation sits in exactly one machine's input queue; a
\textsc{transport} operation's delivery destination is on its transport
record. Matching destinations to machine cells resolves $(\mu,l)$
exactly, so only $r$ need be bounded.

\paragraph{Bounded releases.}
For predecessor $p$ committed to machine $m=\mu(p)$, the successor's
release bound is
\begin{align}
p\in\textsc{machine\_queue}:\quad
\bar r(o) &\ge \mathrm{avail}(m) + p_m, \label{eq:queue}\\
p\in\textsc{transport}:\quad
\bar r(o) &\ge \Delta_{\mathrm{pen}} + p_m, \label{eq:transport}
\end{align}
where $p_m$ is $p$'s duration on $m$ (max over alternatives if
unassigned), $\mathrm{avail}(m)$ accounts for repair time and the
current operation, and $\Delta_{\mathrm{pen}}$ is the single
\emph{penalty parameter}---a generous travel allowance (default 200
steps for maps up to $32{\times}32$). Bound~\eqref{eq:queue} is exact
when $p$ is next in its FIFO queue and a relaxation otherwise
(work ahead of $p$ is not yet included; \S\ref{sec:sound} shows why the
slack is safe); bound~\eqref{eq:transport} is conservative under
Assumption~A1 below. Material sources use the exactly resolved machine;
client-only configuration keys are stripped from the service payload.

\subsection{Soundness}
\label{sec:sound}
We separate two guarantees. Let the \emph{stack release semantics} be:
\textbf{S1}~release times gate the dispatch of transport requests only;
\textbf{S2}~an AGV's pickup physically blocks at the source until the
material exists; \textbf{S3}~a machine starts an operation only when it
has entered its input queue, which happens only on delivery.

\begin{proposition}[Execution soundness, unconditional]
\label{prop:exec}
Under S1--S3, every conservative projection preserves material
precedence at execution: no operation starts before its predecessor's
material is delivered, for \emph{any} $\bar r\ge 0$ (in particular, for
any $\Delta_{\mathrm{pen}}$).
\end{proposition}
\begin{proof}[Proof sketch]
By S1 an under-estimated $\bar r$ can only dispatch a transfer early;
by S2 that transfer waits at its source until the predecessor
completes; by S3 processing cannot begin before delivery. The failure
mode of a too-small bound is an idle AGV camping at a source
(congestion, under-utilization), never a premature start.
\end{proof}

Proposition~\ref{prop:exec} is what makes the penalty knob safe to
sweep: mis-setting $\Delta_{\mathrm{pen}}$ degrades performance, never
correctness. The second guarantee is about the plan, not the execution:

\begin{proposition}[Plan-time validity, conditional]
\label{prop:plan}
Assume \textbf{A1} (remaining transport time of any committed transfer
$\le\Delta_{\mathrm{pen}}$) and FIFO input queues. Then
$\bar r(o)\ge r(o)$ for all successors, so every start time in the
revised plan is a feasible lower bound: the plan executes without
release-induced stalling beyond the bounds' slack.
\end{proposition}

Under A1 violation (extreme congestion), behavior falls back to
Proposition~\ref{prop:exec}: the plan goes stale but execution stays
safe---a two-level guarantee we exploit when sweeping
$\Delta_{\mathrm{pen}}$ down to 10 steps (\S\ref{sec:experiments}).

\paragraph{Implementation.}
The abstraction instantiates in $\sim$90 lines of client-side encoder
changes in our stack, opt-in via a solver-config flag, leaving legacy
behavior bit-identical. The abstraction itself is stack-independent:
any dispatch-by-release simulator with physically enforced material
precedence (S1--S3) admits the same construction.

\section{Recovery Orchestrator}
On top of unblocked revisions we instantiate the policy space
$\mathcal{O}(\tau,\kappa)$: trigger
$\tau\in\{\texttt{never},\texttt{event},\texttt{periodic-}K,
\texttt{myopic}\}$ and scope
$\kappa\in\{\texttt{full},\texttt{partial}\}$ (scoped repair around a
failed machine), with an escalation ladder (full$\to$partial$\to$defer)
and a minimum replan interval. Policies serving as treatments:
\textsc{greedy-reactive} (rule-based, replans by construction),
\textsc{static} (one-shot CP-SAT, $\tau{=}$never), \textsc{full},
\textsc{partial}, and \textsc{myopic}---a value-aware trigger that, at
each qualifying event, exploits the engine's deterministic
snapshot/replay to roll out both branches for $H$ steps under identical
exogenous events and replans only if
$\hat V(s)=\mathrm{prog}_B-\mathrm{prog}_A>\lambda$, where progress is
completed operations (ties broken by accumulated processing time). This
is the simplest instantiation of the when-to-replan decision policy our
results call for; it requires no learning. One caveat: replay restores
the world RNG, so under stochastic scenarios ($S_3$) the rollouts see
the actual future tape---a mild clairvoyance in that arm; under
deterministic scenarios ($S_0,S_1,S_4$) the rollouts are leak-free by
construction, and E7 headline comparisons hold within both groups. A
reseeded leak-free variant (independent rollout tape, common random
numbers across branches) is left as future work.

\section{Experiments}
\label{sec:experiments}
\textbf{Benchmark.} Brandimarte instances $\times$ map families
(corridor maze / open random) $\times$ fleet sizes, on a full-stack
coupled factory simulator: a discrete-event grid factory with a CP-SAT
scheduling service and a rolling EECBS routing service. Episodes are
subprocess-isolated with hard timeouts; all runs, seeds, and manifests
are recorded. The baseline commitment encoding (\S\ref{sec:diagnosis})
is the system's standard mode; soft commitments are an opt-in flag.
\textbf{Campaign (1209 episodes total):} a 948-episode main factorial
campaign (E1--E4) plus 261 episodes of route-backend sensitivity,
arrival factorials, and value-aware-trigger comparison (E5--E7).
\emph{E1} main factorial: $\{$mk01, mk02, mk05, mk07$\}$ $\times$ maze
$\times$ $K\in\{2,4,6\}$ $\times$ $\{S_0, S_1{=}\texttt{m\_down}@150,
S_4{=}\texttt{m\_down}@80{\times}200, S_3{=}\text{stochastic}\}$ $\times$
4 policies $\times$ 3 seeds.
\emph{E2} repeats a subset on the open map.
\emph{E3} legacy-commitment control (trigger-level activation success).
\emph{E4} penalty sweep $\Delta_{\mathrm{pen}}\in\{50,100,200,400\}$,
plus a $\Delta_{\mathrm{pen}}{=}10$ stress test backing
Proposition~\ref{prop:exec}.
\emph{E5} route-backend sensitivity (EECBS vs Python LaCAM/PIBT).
\emph{E6} arrivals factorial (legacy vs soft commitments;
arrival-revision failures).
\emph{E7} myopic value-aware trigger vs static/full.
We report \emph{trigger-level activation success} (share of triggered
revisions that activate; E1--E3) and \emph{arrival-revision failures}
(failed arrival revisions; E6) as distinct metrics.

\subsection{Results}

\paragraph{R1: Soft commitments fix the mechanism (E3 vs E1/E2).}
Under legacy encoding, revision activation succeeds only
$56/(56{+}50)=53\%$ of trigger opportunities; with soft commitments,
$1125/(1125{+}17)=\mathbf{99\%}$ ($216{+}216$ episodes, $2401$
triggered revisions in total). The mechanism bottleneck is real and the
fix removes it.

\paragraph{R2: With the mechanism fixed, passive robustness dominates
disruption response (E1+E2, 864 episodes).}
Table~\ref{tab:imputed} reports makespans with right-censored episodes
(timeouts or non-completion) imputed at the 4096-step cap, which removes
survivorship bias. \textsc{static} beats \textsc{full}/\textsc{partial}
\emph{in every disruption scenario}: paired win rate
$143{:}12$ (61 ties, $n{=}216$), Wilcoxon
$p<10^{-4}$ on each of $S_1, S_4, S_3$. Disruption regret
(Table~\ref{tab:regret}) tells the same story: one-shot plans finish
within $1.0$--$1.06\times$ of their no-disruption makespan even under
200-step outages, whereas event-triggered revision with our conservative
bounds pays $1.25$--$2.25\times$. Slack absorbs bounded disturbances;
replanning under them only re-imposes conservatism.

\begin{table}[t]
\centering
\small
\begin{tabular}{lrrrr}
\toprule
Scen & greedy & static & full & partial \\
\midrule
$S_0$ & 2136 & \textbf{1117} & 1117 & 1117 \\
$S_1$ & 2115 & \textbf{1003} & 1309 & 1309 \\
$S_4$ & 2201 & \textbf{1017} & 1333 & 1333 \\
$S_3$ & 2464 & \textbf{1618} & 2281 & 2281 \\
\bottomrule
\end{tabular}
\caption{Makespan (mean$\pm$std collapsed to mean; censored episodes
imputed at 4096). \textsc{static} dominates all disruption scenarios.}
\label{tab:imputed}
\end{table}

\begin{table}[t]
\centering
\small
\begin{tabular}{lrrr}
\toprule
Scen & static & full/partial & $n$ \\
\midrule
$S_1$ (outage 40) & $0.99{\pm}0.03$ & $1.25{\pm}0.14$ & 46--48 \\
$S_4$ (outage 200) & $1.06{\pm}0.11$ & $1.26{\pm}0.12$ & 44--48 \\
$S_3$ (stochastic) & $1.03{\pm}0.05$ & $2.25{\pm}1.16$ & 37--39 \\
\bottomrule
\end{tabular}
\caption{Disruption regret $\cmax(S)/\cmax(S_0)$ (finished episodes).}
\label{tab:regret}
\end{table}

\paragraph{R3: Scope equivalence.}
\textsc{full} and \textsc{partial} produced \emph{identical} episodes
on every one of 216 paired runs: under single-machine failures the
affected-successor scope already spans the entire reschedulable set, so
scoped repair and full revision coincide. Scope differentiation requires
concurrent multi-machine failures---outside our scenario set.

\paragraph{R4: The penalty knob is real but locally flat (E4).}
Sweeping $\Delta_{\mathrm{pen}}\in\{50,100,200,400\}$ on
$\{$mk01, mk05$\}{\times}\{S_1,S_4\}$ shows no monotone trend
(best mean at $200$: $899$/$817$ vs $939$--$1067$ elsewhere;
$n{=}4$--$6$ per cell after censoring). A stress run at
$\Delta_{\mathrm{pen}}{=}10$---twenty times below the default and far
under any travel bound---completes all episodes with revisions
activating normally and no execution violations, degrading only in
makespan ($1392$--$3801$ vs $\sim$1400 at the default), exactly as
Proposition~\ref{prop:exec} predicts: early-dispatched transfers wait
at their sources rather than start prematurely. The knob exists and its
response is measurable; adaptive or per-commitment penalties remain
open.

\paragraph{R7: A value-aware trigger improves on both extremes (E7, 108
episodes).}
Under local A$^*$ routing (the regime where the routing layer is
weakest and replanning has the most room to matter), we compare
\textsc{static}, \textsc{full}, and \textsc{myopic} on an identical
factorial $\{$mk01, mk05$\}\times K\in\{4,6\}\times\{S_1,S_4,S_3\}
\times$ 3 seeds. \textbf{\textsc{myopic} wins overall}: censored-imputed
mean $2876$ vs $3018$ (\textsc{static}) and $3466$ (\textsc{full});
censoring $14$ vs $15$/$20$; paired $14{:}9$ ($13$ ties) over static and
$21{:}6$ ($9$ ties) over full. The mechanism is surgical: across $453$
trigger evaluations it replanned only $6$ times ($1.3\%$)---almost
always endorsing the law of R2 (leave the plan alone), yet capturing the
rare states where revision pays (largest gain under $S_4$: $2039$ vs
$2296$). A no-learning, snapshot-based lookahead thus instantiates the
when-to-replan decision policy our diagnosis calls for; in this regime
it improves on both never-replan and always-replan with paired
significance, though the margin is modest against static ($14{:}9$) and
the comparison is limited to local A$^*$ routing.

\paragraph{R5: Arrivals (E6, 81 controlled episodes).}
The regime where slack cannot help is work absent from the plan. On the
pre-fix stack we observed total arrival starvation (throughput
$24.4\to0.24$ jobs/ksteps); the controlled E6 factorial
($\{$mk01, mk02, mk05$\}\times$ cadence $\{10..100\}\times$
$\{$legacy, soft$\}\times$ 3 seeds, plus batch anchors) isolates what
the \emph{commitment encoding} itself contributes on the repaired
engine: legacy revisions still fail $1$--$18$ times per cell whenever
an arrival lands mid-commitment, whereas soft commitments fail
\textbf{never} ($0$ across all 72 streaming episodes). Operationally:
soft guarantees completion everywhere ($10/10$, $15/15$ jobs; the worst
legacy cell drops to $7/10$) and lifts throughput on the largest shop
by $23$--$31\%$ (mk05: $8.1\to10.5$ at cadence 20, $9.6\to12.5$ at 50),
while smaller shops are at parity---arrivals benefit from unblocked
reliability rather than from dramatic rescue.

\paragraph{R6: Route-backend sensitivity (E5, 72 episodes).}
Swapping only the routing backend behind the identical protocol
dramatically reshapes the coupled stack: with the production C++
EECBS service, \textsc{static} completes these cells at
$\approx$400--800 steps; with the reference Python LaCAM server every
metric degrades to near-livelock ($\approx$3900, 4/36 errors), and the
Python PIBT server fails outright (22/36 errors even at a relaxed 3s
planning budget). The rolling closed loop is bottlenecked by
\emph{per-replan latency} of the routing backend---consistent with the
routing layer being the binding constraint of hierarchical
decomposition, and a caution for evaluating learned MAPF backends
inside coupled simulators.

\section{Discussion: When \emph{Not} to Replan}
Our factorial evidence inverts the usual closed-loop pitch. Once the
commitment encoding is fixed and revisions actually execute, the
\emph{correct} decision in every disruption scenario we could construct
is to leave the one-shot plan alone: slack plus release-by-plan
semantics absorbs 200-step outages within 6\%, and conservative
revision consistently pays a $25$--$125\%$ premium for responsiveness
that nothing needs. The regime where replanning is structurally
indispensable is \emph{arrivals}---work absent from the plan that no
amount of slack can absorb. This reframes closed-loop factory
scheduling: the research question is not a better trigger policy but a
\emph{decision policy} for when revision is worth its conservatism,
plus commitment encodings whose penalty adapts (per-commitment travel
bounds, learned penalties). Our unblocked stack makes that question
measurable for the first time.

\section{Conclusion}
In coupled scheduling-and-path-finding stacks, the obstacle to
closed-loop operation is not the replanning policy but the commitment
encoding; a sound conservative relaxation lifts activation from 53\% to
99\% and eliminates revision failures outright. With the pathway
repaired, a 1209-episode factorial study shows that disruption response
should, counterintuitively, stay open-loop: slack dominates
conservatism. Replanning earns its keep through reliability under
arrivals---guaranteed completion and up to $31\%$ throughput on larger
shops---and a no-learning, lookahead value-aware trigger captures the
best of both regimes surgically ($1.3\%$ replan rate, improving on both
extremes in the studied regime). The remaining open problem is adaptive
commitment penalties and richer value estimates, which our unblocked
stack now makes measurable.

\bibliographystyle{plain}
\bibliography{references_icaps_en}

\end{document}
